\documentclass[11pt]{article}
\usepackage{amsmath,amssymb,amsthm,mathtools}
\usepackage[margin=1in]{geometry}
\usepackage{enumitem}
\usepackage{hyperref}

\newtheorem{theorem}{Theorem}
\newtheorem{lemma}{Lemma}
\newtheorem{proposition}{Proposition}
\newtheorem{corollary}{Corollary}
\newtheorem{definition}{Definition}
\newtheorem{remark}{Remark}
\newtheorem{example}{Example}

\let\Ker\relax
\DeclareMathOperator{\Ker}{Ker}
\let\tr\relax
\DeclareMathOperator{\tr}{tr}
\let\Fix\relax
\DeclareMathOperator{\Fix}{Fix}
\let\Cap\relax
\DeclareMathOperator{\Cap}{Cap}

\title{Step 37: Character-Source Frame Theorem for Anti-Invariant Anti-Localization}
\author{Six Birds anti-localization working note}
\date{}

\begin{document}
\maketitle

\section{Purpose}
Steps 31--36 isolated a root-composite route to zero confinement.  The anti-linear fixed-locus theorem reduces confinement to two records:
\[
  \mathsf A_Z\preceq \mathsf K^- ,
  \qquad
  \mathsf K^-\preceq \Theta^- .
\]
The first is an explicit-formula domination square.  The second is an anti-invariant anti-localization budget.  Step 36 proved that a strict extension collapses \(\mathsf K^-\) when it supplies full-sector anti-invariant coercivity.

This note specializes that mechanism to the root-composite/character profile.  The main conclusion is:
\[
\boxed{
\text{character-wise source growth + a lower character frame}
\Longrightarrow
\text{full anti-invariant budget collapse}.
}
\]
The key obstruction is also explicit: hardening only invariant, trace, or incomplete character sectors leaves unpriced anti-invariant directions.

\section{Anti-invariant response space and character source records}
Let \(Y^-\) be a finite-dimensional anti-invariant response Hilbert space.  Let \(L^-_\Gamma:E\to Y^-\) be the declared anti-invariant probe family on an exact packaged carrier, and let \(C_\Gamma\succeq0\) be the packaged audit form.  Assume null-mode legality:
\[
  L^-_\Gamma a=0\quad \text{for all }a\in \Ker C_\Gamma.
\]
Then the anti-invariant currency matrix is
\[
  \mathsf K^- = L^-_\Gamma C_\Gamma^\dagger (L^-_\Gamma)^* .
\]

\begin{definition}[character/source record]
A source record is a tuple
\[
  s=(Q_s,\Theta_s,\lambda_s,D_s),
\]
where \(Q_s:Y^-\to Z_s\) is an upstream-visible character or source readout, \(\Theta_s\succ0\) is a local source budget on \(Z_s\), \(\lambda_s>0\) is a declared source strength, and \(D_s\succeq0\) is a positive audit contribution satisfying
\[
  D_s\succeq
  \lambda_s (Q_sL^-_\Gamma)^*\Theta_s^{-1}(Q_sL^-_\Gamma).
\]
The response-side source frame associated to a finite accepted family \(\mathcal S\) is
\[
  F_{\mathcal S}:=
  \sum_{s\in\mathcal S}\lambda_s Q_s^*\Theta_s^{-1}Q_s.
\]
\end{definition}

\begin{definition}[relative lower source-frame constant]
For a global anti-invariant reference budget \(\Theta^-_0\succ0\), define
\[
  \Lambda(\mathcal S;\Theta^-_0)
  :=
  \lambda_{\min}\!\bigl((\Theta^-_0)^{1/2}F_{\mathcal S}(\Theta^-_0)^{1/2}\bigr).
\]
Equivalently,
\[
  F_{\mathcal S}\succeq \Lambda(\mathcal S;\Theta^-_0)(\Theta^-_0)^{-1}.
\]
\end{definition}

\section{The source-frame coercivity theorem}
\begin{theorem}[source-frame anti-invariant coercivity]
Let \(\mathcal S\) be an accepted finite source family.  Suppose the strict extension adds the positive source audit
\[
  C_\Gamma^+ = C_\Gamma + \sum_{s\in\mathcal S}D_s .
\]
If
\[
  \Lambda:=\Lambda(\mathcal S;\Theta^-_0)>0,
\]
then
\[
  C_\Gamma^+\succeq
  \Lambda (L^-_\Gamma)^*(\Theta^-_0)^{-1}L^-_\Gamma,
\]
and therefore
\[
\boxed{
  \mathsf K^-_+ := L^-_\Gamma (C_\Gamma^+)^\dagger (L^-_\Gamma)^*
  \preceq \Lambda^{-1}\Theta^-_0.
}
\]
Consequently, if along an accepted strict-extension ladder \(\Lambda_n\to\infty\), then \(\mathsf K^-_n\to0\) in Loewner order relative to \(\Theta^-_0\).
\end{theorem}

\begin{proof}
By the source-record inequalities,
\[
  \sum_sD_s\succeq
  \sum_s \lambda_s(L^-_\Gamma)^*Q_s^*\Theta_s^{-1}Q_sL^-_\Gamma
  = (L^-_\Gamma)^*F_{\mathcal S}L^-_\Gamma.
\]
If \(F_{\mathcal S}\succeq \Lambda(\Theta^-_0)^{-1}\), then
\[
  C_\Gamma^+\succeq
  (L^-_\Gamma)^*F_{\mathcal S}L^-_\Gamma
  \succeq
  \Lambda (L^-_\Gamma)^*(\Theta^-_0)^{-1}L^-_\Gamma.
\]
For any \(a\in E\), this gives
\[
  \|L^-_\Gamma a\|^2_{(\Theta^-_0)^{-1}}
  \le \Lambda^{-1}\langle a,C_\Gamma^+a\rangle.
\]
Equivalently, for every recombination \(y\in Y^-\),
\[
  |\langle y,L^-_\Gamma a\rangle|^2
  \le
  \Lambda^{-1}\langle y,\Theta^-_0 y\rangle\langle a,C_\Gamma^+a\rangle.
\]
Taking the supremum over legal \(a\) gives
\[
  y^*\mathsf K^-_+y\le \Lambda^{-1}y^*\Theta^-_0y,
\]
which is the claimed Loewner inequality.
\end{proof}

\section{Character decomposition}
Let \(G\) be a finite group acting unitarily on \(Y^-\).  Write the isotypic decomposition
\[
  Y^- = \bigoplus_{\chi\in\widehat G^-}Y_\chi,
  \qquad
  P_\chi:Y^-\to Y_\chi.
\]
Here \(\widehat G^-\) is the set of anti-invariant-visible character blocks.  It need not equal all irreducible characters of \(G\); it is the part of the response family declared native for the anti-invariant ledger.

\begin{corollary}[orthogonal character-source frame]
Assume \(\Theta^-_0\) is block diagonal with respect to \(Y^- =\oplus_\chi Y_\chi\).  For each \(\chi\in\widehat G^-\), let
\[
  Q_\chi=P_\chi,
  \qquad
  \Theta_\chi=P_\chi\Theta^-_0P_\chi|_{Y_\chi},
  \qquad
  \lambda_\chi>0.
\]
Then
\[
  F_{\mathcal S}=\sum_{\chi\in\widehat G^-}\lambda_\chi P_\chi\Theta_\chi^{-1}P_\chi,
\]
and
\[
  \Lambda(\mathcal S;\Theta^-_0)=\min_{\chi\in\widehat G^-}\lambda_\chi.
\]
Thus
\[
  \mathsf K^-_+\preceq
  \bigl(\min_\chi\lambda_\chi\bigr)^{-1}\Theta^-_0.
\]
In particular, character-wise source strengths \(\lambda_\chi(n)\to\infty\) uniformly over \(\widehat G^-\) collapse the anti-invariant budget.
\end{corollary}

\begin{proof}
On each block \(Y_\chi\), the operator \((\Theta^-_0)^{1/2}F_{\mathcal S}(\Theta^-_0)^{1/2}\) acts as \(\lambda_\chi I\).  Hence its smallest eigenvalue is \(\min_\chi\lambda_\chi\).  Apply the theorem.
\end{proof}

\begin{corollary}[nonorthogonal character/source frame]
Suppose \(\Theta_s=\theta_s I_{Z_s}\) and the source readouts obey a lower frame bound
\[
  \sum_{s\in\mathcal S}Q_s^*Q_s\succeq m I_{Y^-}.
\]
If \(\lambda_s/\theta_s\ge c>0\) for every \(s\), then
\[
  F_{\mathcal S}\succeq cm I_{Y^-}.
\]
For \(\Theta^-_0=I\), this gives
\[
  \mathsf K^-_+\preceq (cm)^{-1}I.
\]
\end{corollary}

\section{No-go: incomplete or wrong-sector character sources}
\begin{proposition}[missing character obstruction]
If there exists \(0\ne y\in Y^-\) such that \(Q_sy=0\) for all accepted sources \(s\), then
\[
  \Lambda(\mathcal S;\Theta^-_0)=0.
\]
No finite source-frame coercivity conclusion controls the recombination \(y\); in particular, the strict extension does not collapse the full anti-invariant budget.
\end{proposition}

\begin{proof}
For such \(y\), \(y^*F_{\mathcal S}y=0\).  Since \(\Theta^-_0\succ0\), this forces the relative lower-frame constant to vanish.
\end{proof}

\begin{proposition}[trace/invariant source is insufficient]
A source family supported only on the invariant or trace sector cannot collapse the anti-invariant currency unless the anti-invariant response space is zero or has been lawfully excluded by nonclaim.  Formally, if \(Q_sP_-=0\) for all sources on \(Y^-\), then \(F_{\mathcal S}=0\) on \(Y^-\), so \(\Lambda=0\).
\end{proposition}

\section{RH-facing reading}
For RH the anti-linear involution is
\[
  J(s)=1-\overline{s},
\]
and \(Y^-\) is the anti-invariant response sector measuring displacement from \(\Re(s)=1/2\).  A root-composite character-source theorem would have to identify upstream-visible character/source blocks \(P_\chi\) and source strengths \(\lambda_\chi\) coming from completion, explicit-formula descent, root-composite conjugacy, or threshold/refinement records.  If those blocks form a lower frame over \(Y^-\), and their strengths grow or are large enough, Step 36's coercivity mechanism gives
\[
  \mathsf K^-\preceq \Lambda^{-1}\Theta^-_0.
\]
Together with the completed explicit-formula domination square
\[
  \mathsf A_Z\preceq \mathsf K^- ,
\]
this yields quantitative confinement to the fixed locus.  If \(\Lambda\to\infty\), the visible off-fixed mass is driven to zero.

\begin{remark}[no-smuggling status]
The character/source blocks must be selected by upstream-visible data before the zero-confinement test.  Choosing the blocks because they close
\[
  \mathsf A_Z\preceq\mathsf K^-\preceq\Theta^-
\]
is a target-selected source and has status \texttt{smuggled} or \texttt{support-only}, not accepted.
\end{remark}

\section{Summary}
Step 37 proves that root-composite anti-localization is not merely character decomposition.  It requires a lower character frame: the accepted character/source records must cover the entire anti-invariant native response space.  Uniform source growth over that frame collapses the anti-invariant budget.  Missing characters, invariant-only traces, or wrong-sector hardening leave unpriced anti-invariant witnesses.

\end{document}
